O Quadro~\ref{tab:equiv_nucleo} apresenta a correspondência entre as funções do núcleo matemático do sistema Java legado (\texttt{BaggingOptimized.java}) e as do módulo \texttt{models.py}. O Quadro~\ref{tab:equiv_interface} reúne as funções de interface, entrada e saída de dados, que na reescrita em Python/Streamlit se distribuem entre \texttt{models.py}, \texttt{components.py} e \texttt{app.py}.

\begingroup
\footnotesize
\renewcommand{\arraystretch}{1.15}
\setlength{\tabcolsep}{4pt}

\refstepcounter{quadro}
\addcontentsline{loq}{quadro}{\protect\numberline{\thequadro}Equivalência entre as funções do núcleo matemático: Java e Python}
\begin{center}
  \begin{minipage}{0.85\linewidth}
    \centering
    \quadroname~\thequadro~--~Equivalência entre as funções do núcleo matemático: Java e Python.
  \end{minipage}
\end{center}
\label{tab:equiv_nucleo}%tab:equiv_nucleo
\vspace{-0.3cm}

\begin{longtable}{|>{\raggedright\arraybackslash}p{0.46\textwidth}|>{\raggedright\arraybackslash}p{0.46\textwidth}|}
\hline
\textbf{Função Java (\texttt{BaggingOptimized.java})} & \textbf{Função Python equivalente} \\ \hline
\endfirsthead

\hline
\textbf{Função Java (\texttt{BaggingOptimized.java})} & \textbf{Função Python equivalente} \\ \hline
\endhead

\hline
\multicolumn{2}{r}{\footnotesize\textit{Continua na próxima página}} \\
\endfoot

\hline
\endlastfoot

% ---------- Núcleo matemático (models.py) ----------
\texttt{performRLM(Xcal, Ycal, Xteste)} & \texttt{perform\_rlm(Xcal, Ycal, Xteste)} em \texttt{models.py} \\ \hline
\texttt{performBagging(Xcal, Ycal, Xteste, numBags, withReplacement, sampleFraction)} & \texttt{perform\_bagging(Xcal, Ycal, Xteste, num\_bags, with\_replacement, fraction, seed)} \\ \hline
\texttt{pseudoInverse(matrix)} & \texttt{np.linalg.lstsq(...)}, invocada em \texttt{perform\_rlm()} e \texttt{perform\_bagging()} \\ \hline
\texttt{calculateMetrics(actual, predicted, nanoTime)} & \texttt{calc\_metrics(actual, predicted)}; o tempo é medido em \texttt{run\_full\_analysis()} com \texttt{time.perf\_counter()} \\ \hline
\texttt{calculateRMSE(actual, predicted)} & \texttt{calc\_metrics()} (retorno \texttt{rmse}) e cálculo direto em \texttt{run\_cross\_validation()} \\ \hline
\texttt{calculateMAE(actual, predicted)} & \texttt{calc\_metrics()} (retorno \texttt{mae}) \\ \hline
\texttt{calculateR2(actual, predicted)} & \texttt{calc\_metrics()} (retorno \texttt{r2}) \\ \hline
\texttt{calculateStdDev(values, mean)} & \texttt{np.std()}, invocada em \texttt{run\_cross\_validation()} \\ \hline
\texttt{optimizeBaggingParameters(Xcal, Ycal, m)} & \texttt{run\_cross\_validation(Xcal, Ycal, m, True, 1.0)} \\ \hline
\texttt{optimizeSubaggingParameters(Xcal, Ycal, m, k)} & \texttt{run\_cross\_validation(Xcal, Ycal, m, False, k)} \\ \hline
\texttt{findOptimalParameters(Xcal, Ycal, mValues, kValues, log)} & \texttt{find\_optimal\_parameters(Xcal, Ycal, m\_values, k\_values, progress\_callback, lang)} \\ \hline
\texttt{executarAnalise()} (Etapas 1 a 3 no \texttt{SwingWorker}) & \texttt{run\_full\_analysis(Xcal, Ycal, Xteste, Yteste, m\_values, k\_values, progress\_cb, lang)} \\ \hline
\texttt{extractRows(matrix, indices)} & Indexação NumPy (\texttt{Xcal[train\_idx]}, \texttt{Xcal[val\_idx]}) em \texttt{run\_cross\_validation()} \\ \hline
\texttt{copyRow(src, srcRow, dest, destRow)} & Indexação NumPy (\texttt{Xcal[idx]}) em \texttt{perform\_bagging()} \\ \hline
\texttt{Collections.shuffle(indices, new Random(42))} & \texttt{np.random.default\_rng(seed).permutation(n)} em \texttt{run\_cross\_validation()} \\ \hline
Classes \texttt{CVResult}, \texttt{OptimalParams} e \texttt{FinalMetrics} & Tuplas e dicionários (\texttt{optimal}, \texttt{metrics}), sem classes dedicadas \\ \hline

% ---------- Leitura e validação de dados ----------
\texttt{loadMatrixFromFile(filename, rows, cols)} & \texttt{parse\_matrix(file\_obj)} em \texttt{models.py}; lê \texttt{.txt}, \texttt{.csv} e \texttt{.xlsx}, sem exigir \texttt{rows} e \texttt{cols} \\ \hline
\texttt{loadVectorFromFile(filename, rows)} & \texttt{parse\_vector(file\_obj)} em \texttt{models.py}; lê \texttt{.txt}, \texttt{.csv} e \texttt{.xlsx} \\ \hline
\texttt{detectarDimensoes(file)} & Atributo \texttt{shape} do \texttt{np.ndarray} retornado por \texttt{parse\_matrix()}, registrado no \textit{log} por \texttt{run\_full\_analysis()} \\ \hline
\texttt{contarLinhas(file)} & \texttt{Ycal.shape[0]} e \texttt{Yteste.shape[0]} do vetor retornado por \texttt{parse\_vector()} \\ \hline
\texttt{parseMValues()} & \texttt{parse\_int\_list(text)} em \texttt{app.py} \\ \hline
\texttt{parseKValues()} & \texttt{parse\_float\_list(text)} em \texttt{app.py} \\ \hline

% ---------- Saída e visualização ----------
\texttt{salvarResultados(conteudo)} & \texttt{st.download\_button} em \texttt{app.py}, que exporta o \textit{log} (\texttt{relatorio.txt}) e as métricas (\texttt{metricas.csv}) sem gravar em disco no servidor \\ \hline
\texttt{plotComparison(yteste, ypred, method)} & \texttt{build\_chart(y\_real, y\_pred, title, pred\_color, lang)} em \texttt{components.py}, complementada por \texttt{build\_scatter\_chart(preds, lang)} \\ \hline

% ---------- Interface ----------
\texttt{initComponents()} & Bloco de layout de \texttt{app.py}: \texttt{st.tabs()}, \texttt{get\_header\_html()} e \texttt{APP\_CSS} (\texttt{components.py}) \\ \hline
\texttt{criarPainelArquivo(label, textField, action)} & Trecho de \texttt{app.py} que combina \texttt{st.markdown()} (rótulo e descrição), \texttt{st.file\_uploader()} e \texttt{st.expander()} para cada arquivo \\ \hline
\texttt{selecionarArquivo(tipo, textField, consumer)} (\texttt{JFileChooser}) & \texttt{st.file\_uploader(type=["txt", "csv", "xlsx"])} em \texttt{app.py} \\ \hline
Campos \texttt{fileXcal}, \texttt{fileYcal}, \texttt{fileXteste} e \texttt{fileYteste} & Variáveis \texttt{xcal\_file}, \texttt{ycal\_file}, \texttt{xteste\_file} e \texttt{yteste\_file} (objetos retornados por \texttt{st.file\_uploader()}) em \texttt{app.py} \\ \hline
\end{longtable}

\vspace{-0.5cm}
\begin{center}
  \begin{minipage}{0.92\linewidth}
    \footnotesize Fonte:~\oautor
  \end{minipage}
\end{center}
\endgroup

% ------------------------------------------------------------
% QUADRO 2 — INTERFACE E ENTRADA/SAÍDA
% ------------------------------------------------------------
\begin{quadro}[ht]
\begin{center}
\footnotesize
\captionsetup{justification=centering, labelsep=endash, width=0.75\linewidth}
\caption{Equivalência entre as funções de interface e de entrada e saída: Java e Python.}
\label{tab:equiv_interface}
\begin{threeparttable}
\renewcommand{\arraystretch}{1.15}
\begin{tabularx}{\textwidth}{|>{\raggedright\arraybackslash}X|>{\raggedright\arraybackslash}X|}
\hline
\textbf{Função Java (\texttt{BaggingOptimized.java})} & \textbf{Função Python equivalente} \\ \hline
\texttt{loadMatrixFromFile(filename, rows, cols)} & \texttt{parse\_matrix(file\_obj)} --- \texttt{models.py} \\ \hline
\texttt{loadVectorFromFile(filename, rows)} & \texttt{parse\_vector(file\_obj)} --- \texttt{models.py} \\ \hline
\texttt{detectarDimensoes(file)} e \texttt{contarLinhas(file)} & Atributo \texttt{.shape} do \textit{array} NumPy, registrado no \textit{log} por \texttt{run\_full\_analysis()} --- \texttt{models.py} \\ \hline
\texttt{parseMValues()} & \texttt{parse\_int\_list(text)} --- \texttt{app.py} \\ \hline
\texttt{parseKValues()} & \texttt{parse\_float\_list(text)} --- \texttt{app.py} \\ \hline
\texttt{salvarResultados(conteudo)} & \texttt{st.download\_button(...)} para o \textit{log} (\texttt{relatorio.txt}) e para as métricas (\texttt{metricas.csv}) --- \texttt{app.py} \\ \hline
\texttt{plotComparison(yteste, ypred, method)} (JFreeChart) & \texttt{build\_chart(y\_real, y\_pred, title, pred\_color, lang)} --- \texttt{components.py} (Plotly) \\ \hline
\texttt{initComponents()}, \texttt{criarPainelArquivo()} e \texttt{selecionarArquivo()} & \texttt{st.file\_uploader}, \texttt{st.text\_input} e \texttt{st.tabs} --- \texttt{app.py} \\ \hline
Campos \texttt{fileXcal}, \texttt{fileYcal}, \texttt{fileXteste}, \texttt{fileYteste} & \texttt{st.session\_state} --- \texttt{app.py} \\ \hline
\multicolumn{1}{|c|}{---} & \texttt{build\_scatter\_chart(preds, lang)} --- gráfico de dispersão Real $\times$ Predito --- \texttt{components.py} \\ \hline
\end{tabularx}
\begin{tablenotes}[flushleft]
\footnotesize
\item Fonte:~\oautor
\item \textit{Nota:} O sinal convencional --- indica que não há função equivalente no sistema Java.
\end{tablenotes}
\end{threeparttable}
\end{center}
\end{quadro}
